\documentclass[10pt,a4paper]{amsart}
\usepackage[margin=19mm]{geometry}
\usepackage{amsmath,amssymb,mathtools}
\usepackage[scr=rsfs,cal=euler]{mathalfa}
\usepackage{xfrac}
\usepackage[unicode,hidelinks,bookmarks=false]{hyperref}
\newcommand{\ie}{i.e.\ }
\newcommand{\cf}{cf.\ }
\newcommand{\resp}{respectively\ }
\newcommand{\Subsection}{\subsection}
\providecommand{\nobreakdash}{\nobreak\hbox{-}}
\setlength{\emergencystretch}{2em}
\multlinegap0pt
\usepackage{bm}
\usepackage{bbm}
\usepackage{dsfont}
\usepackage{xspace}
\usepackage{cancel}
\usepackage{mathtools}
\usepackage{amsthm}
\makeatletter
\@ifundefined{thm}{\newtheorem{thm}{Thm}}{}
\@ifundefined{theorem}{\newtheorem{theorem}[thm]{Theorem}}{}
\makeatother
\title[Weighted error-sum identities for periodic continued fractio]{Weighted error-sum identities for periodic continued fractions and their generalizations}
\author{Kevin Calderon, Nikita Kalinin}
\date{Source: arXiv:2601.07862v2 (CC BY 4.0)}
\begin{document}
\maketitle

\noindent\emph{Source and licence.} Weighted error-sum identities for periodic continued fractions and their generalizations. Version arXiv:2601.07862v2 (CC BY 4.0), distributed under the Creative Commons Attribution 4.0 International licence, as recorded in the arXiv licence metadata for this exact version and shown on the abstract page https://arxiv.org/abs/2601.07862v2. Full rights notice: licenses/arxiv-2601.07862-CC-BY-4.0.txt.\par\smallskip

\noindent\textit{Editorial scope.} Counted unit: the Theorem whose source label is \texttt{general} in the article (source0.tex lines 1435--1450, source label \texttt{general}), delivered together with the article's own complete original proof of it (source0.tex lines 1452--1486, one \texttt{proof} environment of 35 lines). No mathematics is written, repaired or re-derived: statement and proof are the article's own text line for line. Required context retained uncounted: none. Every definition, display and label of those lines is retained unchanged. Omitted from this package and not redistributed: 1--1434, 1451--1451, 1487--1642. Nothing inside a retained excerpt is omitted or altered: every retained body is delivered exactly as the article's lines are written, so no omission receipt of any kind accompanies this package. Citations: the retained excerpts carry no citation command at all, so no bibliography entry of the article is transcribed and no reference is redistributed. Reference adaptation: none was needed and none was made; every reference inside the delivered unit resolves inside it, so no unresolved reference or citation is produced. Printed slips kept literal: no printed word, symbol, hypothesis or number of the counted statement or of its proof was altered. Layout and class: the article's own class is \texttt{interact} with packages including latexsym, anyfontsize, tikz, color, hyperref, url, breakurl; the wrapper is the lane's pinned \texttt{amsart} editorial wrapper at 10pt on A4, which restates the theorem-like environments the retained text needs and adds the article's own macro definitions that the retained text needs (none), with extra wrapper packages (bm, bbm, dsfont, xspace, cancel, mathtools). The delivered numbering is the unit's own. Assets: no retained span contains a figure environment, a graphics-inclusion command, a PDF-page inclusion, a table or a TikZ picture; no image file is copied into this package, the package declares no asset dependency and no image file is redistributed with it. Licence: version arXiv:2601.07862v2 of this article is distributed under the Creative Commons Attribution 4.0 International licence, as recorded in the arXiv licence metadata for this exact version and shown on the abstract page https://arxiv.org/abs/2601.07862v2. A full rights notice is delivered at licenses/arxiv-2601.07862-CC-BY-4.0.txt. Characters: every retained excerpt is pure ASCII, so the delivered engine, XeLaTeX with its default Latin Modern text font, reports no missing character.
\begin{theorem}\label{general}

Suppose that the JPA expansion of \(\eta\) is purely periodic and that, in the
chosen convention, the corresponding convergent columns satisfy
\[
C_{mN+r}=M_N^mC_r.
\]
Assume moreover that
\[
M_N^T\mathcal B=\mu\mathcal B
\]
for some \(\mu\) with \(|\mu|<1\). Then for any \(m\ge0\) and residue \(r\) we have

 $$\vartheta_{mN+r}=\mu^m \vartheta_{r}.$$
    
\end{theorem}

\begin{proof}
From the definition of the algorithm being purely periodic with period \(N\), the
convergent vector at step \(mN+r\) is
\[
C_{mN+r}=M_N^mC_r.
\]
Therefore
\[
\vartheta_{mN+r}
=
\langle M_N^mC_r,\mathcal B\rangle.
\]
Using
\[
\langle Ax,y\rangle=\langle x,A^Ty\rangle,
\]
we get
\[
\vartheta_{mN+r}
=
\langle C_r,(M_N^T)^m\mathcal B\rangle.
\]
By hypothesis,
\[
(M_N^T)^m\mathcal B=\mu^m\mathcal B.
\]
Hence
\[
\vartheta_{mN+r}
=
\mu^m\langle C_r,\mathcal B\rangle
=
\mu^m\vartheta_r.
\]
\end{proof}

\end{document}
